% META: {"id": "1SPE-PROBCOND-EX-023", "chapitre": "1SPE-PROBA-COND", "type_objet": "exercice", "capacites": ["1SPE-PROBA-COND-C3"], "capacites_codes": ["C3"], "methodes": ["M3"], "parcours": 1, "competences": ["calculer", "raisonner"], "duree_min": 12, "mode_creation": "ex_nihilo", "sources_inspiration": [], "parametres_sympy": {}, "fichier_tex": "chapitres/1SPE-PROBA-COND/exercices/1SPE-PROBCOND-EX-023.tex", "corrige_tex": "chapitres/1SPE-PROBA-COND/corriges/1SPE-PROBCOND-CO-023.tex", "status": "generated"}
% BEGIN-VERIFY
% from sympy import Rational
% P_C1 = Rational(1, 2)
% P_C2 = Rational(3, 10)
% P_C3 = Rational(1, 5)
% assert P_C1 + P_C2 + P_C3 == 1
% P_D_C1 = Rational(1, 50)
% P_D_C2 = Rational(3, 100)
% P_D_C3 = Rational(1, 20)
% P_D = P_C1*P_D_C1 + P_C2*P_D_C2 + P_C3*P_D_C3
% assert P_D == Rational(1,100) + Rational(9,1000) + Rational(1,100)
% assert P_D == Rational(10,1000) + Rational(9,1000) + Rational(10,1000)
% assert P_D == Rational(29, 1000)
% END-VERIFY
\begin{exercice}{1SPE-PROBCOND-EX-023}{1}{12}
Une usine dispose de trois chaines : $C_1$ ($50\,\%$), $C_2$ ($30\,\%$), $C_3$ ($20\,\%$). Les taux de defaut sont respectivement $2\,\%$, $3\,\%$ et $5\,\%$.

\begin{enumerate}
  \item Verifier que $\{C_1, C_2, C_3\}$ est une partition.
  \item Calculer $P(D)$ ou $D$ est << piece defectueuse >>.
\end{enumerate}
\end{exercice}
